\section{Independent-Seed Control}
\label{app:seed-control}

\begingroup\setlength{\parskip}{2pt plus 1pt}

To test for topic-level leakage from shared seeds across generators, we construct an independent-seed evaluation where Pythia-2.8B and OLMo-2-1B use completely disjoint Wikipedia seed pools (500 seeds each, zero overlap).
Each generator produces depth-0 through depth-4 chains on its own seed pool; features are extracted with the same Pythia-410M reference model.
We then train a 7-feature XGBoost classifier on one generator's data and test on the other, in both directions.

\begin{table}[ht]
\centering
\small
\caption{Independent-seed cross-generator evaluation (binary: gen-1 vs.\ gen-2+). Disjoint seed pools eliminate any possibility of topic-based leakage. The average independent-seed F1 (0.680) is within 1.4\,pp of the shared-seed LOGO baseline (0.694). The classifier thus captures generative degradation patterns rather than seed-specific topic content.}
\label{tab:seed-control}
\resizebox{\linewidth}{!}{%
\renewcommand{\arraystretch}{1.15}%
\begin{tabular}{lcc}
\toprule
\textbf{Setting} & \textbf{Binary F1} & \textbf{95\% CI} \\
\midrule
Pythia-2.8B $\rightarrow$ OLMo-2-1B & 0.695 & [0.673, 0.717] \\
OLMo-2-1B $\rightarrow$ Pythia-2.8B & 0.666 & [0.643, 0.689] \\
\midrule
Independent-seed average & 0.680 & --- \\
Shared-seed LOGO (reference) & 0.694 & [0.688, 0.700] \\
\bottomrule
\end{tabular}}
\end{table}

The small gap ($-$1.4\,pp) between independent-seed and shared-seed evaluation falls within bootstrap confidence intervals, and both directions individually exceed 0.66.
Combined with the seed-level demeaning control (\S\ref{sec:results-classification}), these results provide converging evidence that depth classification is driven by distributional properties of recursive generation, not topic identity.

\paragraph{Semantic similarity analysis.}
Embedding cosine similarity between consecutive depth steps increases monotonically (Figure~\ref{fig:semantic-sim}): 0$\rightarrow$1 (0.622), 1$\rightarrow$2 (0.683), 2$\rightarrow$3 (0.801), 3$\rightarrow$4 (0.859).
Deeper generations are thus \emph{more} semantically similar to their inputs, yet the classifier still distinguishes depth classes. The depth signal resides in distributional rather than topical properties.
Prompt type modulates the baseline similarity level (continuation 0.717 $<$ elaboration 0.809 $\approx$ rewriting 0.808), plausibly because semantic-preservation constraints reduce topic drift while leaving the distributional location shift intact.

\begin{figure}[t]
\centering
\includegraphics[width=0.9\columnwidth]{fig_appendix_semantic_sim.pdf}
\caption{Embedding cosine similarity between consecutive depth steps. Deeper generations become increasingly similar to their inputs, yet the classifier still distinguishes depth classes. Dotted lines: prompt-type baselines.}
\label{fig:semantic-sim}
\end{figure}

\endgroup
